\chapter{一元函数微分学}

\section{原题页索引}
本章原题对应做题本第 19--38 页。原题页图片已随 Overleaf 项目打包。

\begin{center}
\includegraphics[width=.92\textwidth]{figures/original_pages/page_019.png}
\end{center}

\section{选择题}

\begin{problemblock}
\textbf{1.} 设 \(f(x)\) 在 \(x=0\) 处连续，则 \(f(x)\) 在 \(x=0\) 处可导的充分条件是（\quad）

A. \(\displaystyle\lim_{x\to0}\frac{f(x)-f(-x)}{2x}\) 存在。

B. \(\displaystyle\lim_{x\to0}\frac{f(\ln(1+x^2))-f(0)}{x^2}\) 存在。

C. \(\displaystyle\lim_{x\to0}\frac{f(x)-f(0)}{\sqrt[3]{x}}\) 存在。

D. \(\displaystyle\lim_{x\to\infty}x f\left(\frac1x\right)\) 存在。
\end{problemblock}

\begin{solutionblock}
\analysis{题目问“充分条件”，要找能推出导数定义极限存在的选项。}
对 D，令 \(t=\frac1x\)，则 \(x\to\infty\) 等价于 \(t\to0\)，且
\[
x f\left(\frac1x\right)=\frac{f(t)}{t}.
\]
若该极限存在且有限，由 \(f\) 在 0 处连续可知 \(f(t)\to f(0)\)。若 \(\frac{f(t)}t\) 有有限极限，则必有 \(f(0)=0\)，否则 \(\frac{f(t)}t\) 会发散。因此
\[
f'(0)=\lim_{t\to0}\frac{f(t)-f(0)}t
=\lim_{t\to0}\frac{f(t)}t
\]
存在，故 D 是充分条件。

其余选项不能推出可导。A 只控制对称差商，例如 \(f(x)=|x|\) 时 A 中极限为 0，但 \(f\) 在 0 不可导。B 只沿 \(\ln(1+x^2)\ge0\) 的右侧路径控制变化，不能保证双侧导数。C 中分母是 \(\sqrt[3]x\)，即使极限存在，也不能推出 \(\frac{f(x)-f(0)}x\) 存在。

\answer{D}
\examnote{可导的核心是 \(\frac{f(x)-f(0)}x\) 的双侧极限；偏路径、对称差商、换错尺度都不够。}
\end{solutionblock}

\begin{problemblock}
\textbf{2.} 设
\[
f(x)=
\begin{cases}
x^2\sin\frac1x,&x\ne0,\\
0,&x=0,
\end{cases}
\]
则在点 \(x=0\) 处函数 \(f(x)\)（\quad）

A. 不连续 \qquad B. 连续但不可导

C. 可导但导数不连续 \qquad D. 可导但导数连续
\end{problemblock}

\begin{solutionblock}
\analysis{这是经典振荡函数题，分别检查连续性、可导性和导函数连续性。}
连续性：
\[
|x^2\sin(1/x)|\le x^2\to0,
\]
故 \(f(x)\to f(0)=0\)，函数在 0 连续。

可导性：
\[
f'(0)=\lim_{x\to0}\frac{f(x)-f(0)}x
=\lim_{x\to0}x\sin\frac1x=0.
\]
故 \(f\) 在 0 可导。

当 \(x\ne0\) 时，
\[
f'(x)=2x\sin\frac1x-\cos\frac1x.
\]
其中 \(-\cos\frac1x\) 在 \(x\to0\) 时振荡不收敛，所以
\[
\lim_{x\to0}f'(x)
\]
不存在。故导数在 0 不连续。

\answer{C}
\examnote{\(x^2\sin(1/x)\) 是“可导但导数不连续”的标准模型。}
\end{solutionblock}

\begin{problemblock}
\textbf{3.} 设函数 \(y=f(x)\) 在点 \(x=0\) 处连续，且
\[
\lim_{x\to0}\frac{f(x)-2x}{1-\cos x}=1,
\]
则 \(f(x)\) 在点 \(x=0\) 处（\quad）

A. 不可导 \qquad B. 可导且 \(f'(0)=0\)

C. 可导且 \(f'(0)=-2\) \qquad D. 可微且 \(\left.dy\right|_{x=0}=2\,dx\)
\end{problemblock}

\begin{solutionblock}
\analysis{由已知极限得到 \(f(x)\) 在 0 附近的主部。}
因为
\[
1-\cos x\sim\frac{x^2}{2},
\]
且
\[
\frac{f(x)-2x}{1-\cos x}\to1,
\]
所以
\[
f(x)-2x\sim1-\cos x\sim\frac{x^2}{2}.
\]
因此
\[
f(x)=2x+O(x^2).
\]
由连续性和上式可得 \(f(0)=0\)。于是
\[
f'(0)=\lim_{x\to0}\frac{f(x)-f(0)}x
=\lim_{x\to0}\frac{2x+O(x^2)}x=2.
\]
一元函数在点处可导即在该点可微，因此
\[
\left.dy\right|_{x=0}=f'(0)\,dx=2\,dx.
\]

\answer{D}
\examnote{已知 \(f(x)-2x\) 是二阶小量，就能立刻读出一阶主部 \(f(x)\sim2x\)。}
\end{solutionblock}

\begin{problemblock}
\textbf{4.} 若 \(f(x)\) 在点 \(x_0\) 处的左、右导数都存在，则 \(f(x)\) 在点 \(x_0\) 处（\quad）

A. 可导 \qquad B. 连续 \qquad C. 不可导 \qquad D. 不一定连续
\end{problemblock}

\begin{solutionblock}
\analysis{单侧导数存在会推出对应单侧连续，但左右导数未必相等。}
若左导数存在，则
\[
\lim_{x\to x_0^-}\frac{f(x)-f(x_0)}{x-x_0}
\]
存在，从而
\[
f(x)-f(x_0)=(x-x_0)\cdot O(1)\to0\quad(x\to x_0^-).
\]
所以左连续。同理，右导数存在推出右连续。左右都连续，故 \(f\) 在 \(x_0\) 连续。

但左右导数存在不代表相等，例如 \(f(x)=|x|\) 在 0 处左导数为 \(-1\)，右导数为 \(1\)，不可导。因此只能推出连续。

\answer{B}
\examnote{“左右导数都存在”比“连续”强，但比“可导”弱；还差左右导数相等。}
\end{solutionblock}

\begin{problemblock}
\textbf{5.} 已知 \(f(x)\) 在 \(x=0\) 处连续，且
\[
\lim_{x\to0}\left[f(x)+e^x\right]^{1/x}=2,
\]
则 \(f'(0)\)（\quad）

A. 不存在 \qquad B. 等于 \(\ln2\) \qquad C. 等于 \(2\) \qquad D. 等于 \(\ln2-1\)
\end{problemblock}

\begin{solutionblock}
\analysis{这是 \(1^\infty\) 型极限反求函数一阶主部。}
设
\[
g(x)=f(x)+e^x.
\]
由
\[
g(x)^{1/x}\to2
\]
可知 \(g(x)\to1\)。又 \(f\) 在 0 连续，故
\[
f(0)+1=1,\qquad f(0)=0.
\]
取对数：
\[
\lim_{x\to0}\frac{\ln g(x)}x=\ln2.
\]
由于 \(g(x)\to1\)，有
\[
\ln g(x)\sim g(x)-1.
\]
所以
\[
\lim_{x\to0}\frac{g(x)-1}{x}=\ln2.
\]
而
\[
g(x)-1=f(x)+e^x-1=f(x)+x+o(x).
\]
于是
\[
\lim_{x\to0}\frac{f(x)}x+1=\ln2,
\]
故
\[
f'(0)=\ln2-1.
\]

\answer{D}
\examnote{先由极限推出底数趋于 1，再取对数，是处理 \(1^\infty\) 的固定动作。}
\end{solutionblock}

\begin{problemblock}
\textbf{6.} 设 \(f(x)\) 有连续一阶导数，\(f(0)=0\)，若当 \(x\to0\) 时，
\[
\int_0^{f(x)}f(t)\,dt
\]
与 \(4x^2\) 为等价无穷小，则 \(f'(0)\) 等于（\quad）

A. 0 \qquad B. 1 \qquad C. 2 \qquad D. \(\frac12\)
\end{problemblock}

\begin{solutionblock}
\analysis{令 \(f'(0)=a\)，把 \(f(x)\) 与 \(f(t)\) 都用一阶主部表示。}
由于 \(f(0)=0\) 且 \(f\) 可导，
\[
f(x)=ax+o(x).
\]
同理
\[
f(t)=at+o(t).
\]
于是
\[
\int_0^{f(x)}f(t)\,dt
\sim\int_0^{ax}at\,dt
=\frac{a}{2}(ax)^2
=\frac{a^3}{2}x^2.
\]
它与 \(4x^2\) 等价，故
\[
\frac{a^3}{2}=4,\qquad a^3=8,\qquad a=2.
\]

\answer{C}
\examnote{变上限也是 \(f(x)\)，所以阶数和系数都会出现 \(a\)，最后是 \(a^3\)。}
\end{solutionblock}

\begin{problemblock}
\textbf{7.} 函数
\[
f(x)=|x-x^2|(e^x-1)+\sin|x-2|
\]
不可导点的个数为（\quad）

A. 0 \qquad B. 1 \qquad C. 2 \qquad D. 3
\end{problemblock}

\begin{solutionblock}
\analysis{不可导主要来自绝对值尖点，但若绝对值因子被零因子抵消，也可能变得可导。}
\(|x-x^2|=|x(1-x)|\) 的尖点候选为 \(x=0,1\)。在 \(x=0\) 附近，
\[
|x-x^2|(e^x-1)\sim |x|\cdot x=x|x|,
\]
其在 0 处导数左右都为 0，故 \(x=0\) 可导。在 \(x=1\) 处，\(e^x-1\ne0\)，绝对值尖点不能被抵消，所以不可导。

\(\sin|x-2|\) 在 \(x=2\) 处不可导，因为
\[
\sin|x-2|\sim |x-2|.
\]
其他点由初等函数复合可导。因此不可导点为 \(x=1,2\)，共 2 个。

\answer{C}
\examnote{绝对值零点是候选点，但要检查外面的乘子是否把尖点阶数抹平。}
\end{solutionblock}

\begin{problemblock}
\textbf{8.}
\[
f(x)=\lim_{n\to\infty}\sqrt[n]{1+|x|^n+e^{nx}},
\]
不可导点的个数为（\quad）

A. 0 \qquad B. 1 \qquad C. 2 \qquad D. 3
\end{problemblock}

\begin{solutionblock}
\analysis{利用结论 \(\lim_{n\to\infty}(a^n+b^n+c^n)^{1/n}=\max\{a,b,c\}\)。}
这里
\[
f(x)=\max\{1,|x|,e^x\}.
\]
当 \(x<-1\) 时，\(|x|>1\) 且 \(e^x<1\)，故 \(f(x)=|x|=-x\)。
当 \(-1\le x\le0\) 时，\(f(x)=1\)。
当 \(x>0\) 时，\(e^x>1\) 且 \(e^x>x\)，故 \(f(x)=e^x\)。

因此
\[
f(x)=
\begin{cases}
-x,&x<-1,\\
1,&-1\le x\le0,\\
e^x,&x>0.
\end{cases}
\]
在 \(x=-1\) 处左右导数分别为 \(-1,0\)，不可导；在 \(x=0\) 处左右导数分别为 \(0,1\)，不可导。其余各段可导。

\answer{C}
\examnote{\(n\) 次根极限通常转化为“最大值函数”，不可导点出现在最大项切换处。}
\end{solutionblock}

\begin{problemblock}
\textbf{9.} 已知 \(f(x)\) 在 \(x=0\) 处连续，且
\[
\lim_{x\to0}\frac{x^2}{f(x)}=1,
\]
则下列结论中正确的个数为（\quad）

① \(f'(0)\) 存在，且 \(f'(0)=0\)。

② \(f''(0)\) 存在，且 \(f''(0)=2\)。

③ \(f(x)\) 在 \(x=0\) 处取得极小值。

④ \(f(x)\) 在 \(x=0\) 的某邻域内连续。

A. 1 \qquad B. 2 \qquad C. 3 \qquad D. 4
\end{problemblock}

\begin{solutionblock}
\analysis{已知条件等价于 \(f(x)\sim x^2\)。它能给出函数值和一阶导数信息，但不能保证二阶导数或邻域连续。}
由 \(f(x)\sim x^2\) 且 \(f\) 在 0 连续，得
\[
f(0)=0.
\]
于是
\[
f'(0)=\lim_{x\to0}\frac{f(x)-f(0)}x
=\lim_{x\to0}\frac{f(x)}x
=\lim_{x\to0}x\cdot\frac{f(x)}{x^2}=0,
\]
故①正确。

②不一定正确。例如可令
\[
f(x)=x^2+x^3\sin\frac1{x^2}\quad(x\ne0),\qquad f(0)=0,
\]
仍有 \(f(x)\sim x^2\)，但二阶导数在 0 处不一定存在。

③正确。因为 \(f(x)\sim x^2>0\)，所以存在去心邻域使 \(f(x)>0=f(0)\)，故 0 处取得极小值。

④不一定正确。极限 \(f(x)\sim x^2\) 只保证 \(f\) 在 0 处的趋近行为，不能保证它在 0 的整个某邻域内每一点都连续。

因此正确的是①③，共 2 个。

\answer{B}
\examnote{“等价于 \(x^2\)”能推出极小值和一阶导数，但不能自动推出二阶可导。}
\end{solutionblock}

\begin{problemblock}
\textbf{10.} 设函数 \(f(x)\) 在 \((-\infty,+\infty)\) 内连续，其导函数的图形如原题图所示，则 \(f(x)\) 有（\quad）

A. 一个极小值点和两个极大值点。

B. 两个极小值点和一个极大值点。

C. 两个极小值点和两个极大值点。

D. 三个极小值点和一个极大值点。
\end{problemblock}

\begin{solutionblock}
\analysis{函数 \(f\) 的极值由导函数 \(f'(x)\) 的符号变化决定：\(f'\) 由正变负，对应极大值；由负变正，对应极小值。}
从图像读出，导函数在四个横轴交点附近的符号变化依次为：
\[
+\to-,
\qquad
-\to+,
\qquad
+\to-,
\qquad
-\to+.
\]
因此第 1 个与第 3 个变号点对应 \(f(x)\) 的极大值点，第 2 个与第 4 个变号点对应 \(f(x)\) 的极小值点。

所以 \(f(x)\) 有两个极小值点和两个极大值点。

\answer{C}
\examnote{看导函数图像时，只看 \(f'\) 在零点两侧的正负号，不要把导函数自己的高低误认为原函数极值。}
\end{solutionblock}

\begin{problemblock}
\textbf{11.} 设函数
\[
f(x)=|x^2(x+1)|
\]
的驻点个数为 \(m\)，极值点的个数为 \(n\)，则（\quad）

A. \(m=1,n=1\) \qquad B. \(m=1,n=2\) \qquad C. \(m=2,n=3\) \qquad D. \(m=3,n=2\)
\end{problemblock}

\begin{solutionblock}
\analysis{先按 \(x=-1\) 分段去绝对值，再找可导点的导数零点与极值点。}
因为 \(x^2\ge0\)，符号由 \(x+1\) 决定：
\[
f(x)=
\begin{cases}
-x^2(x+1),&x<-1,\\
x^2(x+1),&x\ge-1.
\end{cases}
\]
当 \(x<-1\) 时，
\[
f'(x)=-(3x^2+2x)=-x(3x+2),
\]
在该区间没有零点。
当 \(x>-1\) 时，
\[
f'(x)=3x^2+2x=x(3x+2),
\]
零点为
\[
x=0,\qquad x=-\frac23.
\]
二者都是可导点且导数为 0，所以驻点个数
\[
m=2.
\]

极值点方面，\(x=-1\) 处 \(f(x)\ge0\) 且 \(f(-1)=0\)，是极小值点；\(x=-\frac23\) 处导数由正变负，是极大值点；\(x=0\) 处导数由负变正，是极小值点。故极值点共有 3 个：
\[
n=3.
\]

\answer{C}
\examnote{驻点必须是可导且导数为零的点；不可导点也可能是极值点，但不是驻点。}
\end{solutionblock}

\begin{problemblock}
\textbf{12.} 函数
\[
f(x)=\int_{-\pi}^{\pi}(t-x\sin t)^2\,dt
\]
的极值点为（\quad）

A. \(x=2\) 为极小值点 \qquad B. \(x=2\) 为极大值点

C. \(x=1\) 为极小值点 \qquad D. \(x=1\) 为极大值点
\end{problemblock}

\begin{solutionblock}
\analysis{本题的积分变量是 \(t\)，\(x\) 是参数。把积分展开后，\(f(x)\) 是关于 \(x\) 的二次函数。}
展开被积函数：
\[
(t-x\sin t)^2=t^2-2xt\sin t+x^2\sin^2t.
\]
因此
\[
f(x)=\int_{-\pi}^{\pi}t^2\,dt
-2x\int_{-\pi}^{\pi}t\sin t\,dt
x^2\int_{-\pi}^{\pi}\sin^2t\,dt.
\]
其中
\[
\int_{-\pi}^{\pi}t\sin t\,dt=2\int_0^\pi t\sin t\,dt=2\pi,
\]
\[
\int_{-\pi}^{\pi}\sin^2t\,dt=\pi.
\]
所以
\[
f(x)=C-4\pi x+\pi x^2
=\pi(x-2)^2+C-4\pi,
\]
其中 \(C=\int_{-\pi}^{\pi}t^2\,dt\) 为常数。故 \(f(x)\) 在
\[
x=2
\]
处取得极小值。

\answer{A}
\examnote{定积分上下限都是常数时，被积函数中的 \(x\) 要看作参数；展开后常常变成普通函数极值题。}
\end{solutionblock}

\begin{problemblock}
\textbf{13.} 设函数 \(f(x)\) 有二阶导数，且
\[
\lim_{x\to0}\frac{f(x)-a}{\ln(1+x)}=0,
\qquad
\lim_{x\to0}\frac{f''(x)-1}{e^{x^2}-1}=2025,
\]
则（\quad）

A. \(f(0)\) 是 \(f(x)\) 的极大值。

B. \(f(0)\) 是 \(f(x)\) 的极小值。

C. \((0,f(0))\) 是曲线 \(y=f(x)\) 的拐点。

D. \(f(0)\) 不是 \(f(x)\) 的极值，\((0,f(0))\) 也不是曲线 \(y=f(x)\) 的拐点。
\end{problemblock}

\begin{solutionblock}
\analysis{第一个极限给出 \(f(0)\) 和 \(f'(0)\)，第二个极限给出 \(f''(0)\)。}
因为 \(\ln(1+x)\sim x\)，且
\[
\frac{f(x)-a}{\ln(1+x)}\to0,
\]
所以
\[
f(x)-a=o(x).
\]
于是
\[
f(0)=a,\qquad f'(0)=0.
\]
又
\[
e^{x^2}-1\sim x^2,
\]
第二个极限说明
\[
f''(x)-1=2025x^2+o(x^2),
\]
因此
\[
f''(0)=1>0.
\]
由二阶导数判别法，\(f(0)\) 是极小值。

\answer{B}
\examnote{先从 \(f(x)-a=o(x)\) 读出 \(f(0)=a,f'(0)=0\)，再用二阶导数判极值。}
\end{solutionblock}

\begin{problemblock}
\textbf{14.} 设函数 \(f(x)\) 有二阶连续导数，且
\[
f(0)=0,\qquad f'(0)>0,\qquad f''(0)<0,
\]
则（\quad）

A. \(x=0\) 是 \(|f(x)|\) 的极值点，但 \((0,f(0))\) 不是曲线 \(y=|f(x)|\) 的拐点。

B. \(x=0\) 不是 \(|f(x)|\) 的极值点，但 \((0,f(0))\) 是曲线 \(y=|f(x)|\) 的拐点。

C. \(x=0\) 是 \(|f(x)|\) 的极值点，且 \((0,f(0))\) 是曲线 \(y=|f(x)|\) 的拐点。

D. \(x=0\) 不是 \(|f(x)|\) 的极值点，且 \((0,f(0))\) 不是曲线 \(y=|f(x)|\) 的拐点。
\end{problemblock}

\begin{solutionblock}
\analysis{由于 \(f(0)=0,f'(0)>0\)，函数 \(f(x)\) 在 0 附近与 \(x\) 同号。}
当 \(x<0\) 且充分接近 0 时，\(f(x)<0\)，故
\[
|f(x)|=-f(x).
\]
当 \(x>0\) 且充分接近 0 时，\(f(x)>0\)，故
\[
|f(x)|=f(x).
\]
因此 \(|f(x)|\ge0=|f(0)|\)，所以 \(x=0\) 是 \(|f(x)|\) 的极小值点。

再看凹凸性：左侧
\[
(|f|)''=(-f)''=-f''.
\]
由于 \(f''(0)<0\)，左侧附近 \(-f''(x)>0\)，曲线凹向上；右侧
\[
(|f|)''=f''(x)<0,
\]
曲线凹向下。凹凸性在 0 两侧发生变化，因此 \((0,f(0))=(0,0)\) 是曲线 \(y=|f(x)|\) 的拐点。

\answer{C}
\examnote{\(|f(x)|\) 在 \(f(x)=0\) 且 \(f'(0)\ne0\) 处形成尖点，同时也可能发生凹凸性改变。}
\end{solutionblock}

\begin{problemblock}
\textbf{15.} 设 \(f(x)\) 满足
\[
f'(0)=0,\qquad f'(x)+[f(x)]^3=x^2,
\]
则（\quad）

A. \(f(0)\) 是 \(f(x)\) 的极大值。

B. \(f(0)\) 是 \(f(x)\) 的极小值。

C. \((0,f(0))\) 是曲线 \(y=f(x)\) 的拐点。

D. \(f(0)\) 不是 \(f(x)\) 的极值，\((0,f(0))\) 也不是曲线 \(y=f(x)\) 的拐点。
\end{problemblock}

\begin{solutionblock}
\analysis{由方程逐次代入和求导，求出 \(f\) 在 0 处的低阶导数。}
令 \(x=0\)，得
\[
f'(0)+[f(0)]^3=0.
\]
又 \(f'(0)=0\)，所以
\[
f(0)=0.
\]
对
\[
f'(x)+[f(x)]^3=x^2
\]
求导：
\[
f''(x)+3[f(x)]^2f'(x)=2x.
\]
令 \(x=0\)，得
\[
f''(0)=0.
\]
再求导：
\[
f'''(x)+6f(x)[f'(x)]^2+3[f(x)]^2f''(x)=2.
\]
令 \(x=0\)，得
\[
f'''(0)=2.
\]
因此
\[
f(x)=\frac{f'''(0)}{6}x^3+o(x^3)=\frac{x^3}{3}+o(x^3).
\]
这说明 0 附近函数像三次函数一样穿过原点，不是极值点；且
\[
f''(x)=2x+o(x)
\]
在 0 两侧变号，因此 \((0,0)\) 是拐点。

\answer{C}
\examnote{当一阶、二阶导数都为 0 时，继续看三阶主部；三次主部通常对应拐点而非极值。}
\end{solutionblock}

\begin{problemblock}
\textbf{16.} 曲线
\[
y=\frac{x^2+1}{\sqrt{x^2-1}}
\]
的渐近线条数为（\quad）

A. 1 \qquad B. 2 \qquad C. 3 \qquad D. 4
\end{problemblock}

\begin{solutionblock}
\analysis{定义域为 \(|x|>1\)，先找垂直渐近线，再看无穷远处斜渐近线。}
当 \(x\to1^+\) 或 \(x\to-1^-\) 时，\(\sqrt{x^2-1}\to0^+\)，分子趋于 2，故有两条垂直渐近线：
\[
x=1,\qquad x=-1.
\]
当 \(x\to+\infty\) 时，
\[
\frac{x^2+1}{\sqrt{x^2-1}}-x\to0,
\]
故有斜渐近线
\[
y=x.
\]
当 \(x\to-\infty\) 时，\(\sqrt{x^2-1}\sim |x|=-x\)，
\[
\frac{x^2+1}{\sqrt{x^2-1}}+x\to0,
\]
故有斜渐近线
\[
y=-x.
\]
共有 4 条渐近线。

\answer{D}
\examnote{含 \(\sqrt{x^2}\) 的无穷远判断要分 \(+\infty\) 与 \(-\infty\)，因为 \(\sqrt{x^2}=|x|\)。}
\end{solutionblock}

\begin{problemblock}
\textbf{17.} 曲线
\[
y=\frac{x^2+x}{x^2-1}
\]
的渐近线条数为（\quad）

A. 0 \qquad B. 1 \qquad C. 2 \qquad D. 3
\end{problemblock}

\begin{solutionblock}
\analysis{先因式分解，区分可去间断点和垂直渐近线。}
\[
y=\frac{x(x+1)}{(x-1)(x+1)}=\frac{x}{x-1}\quad(x\ne\pm1).
\]
在 \(x=-1\) 处只是可去间断点，不是渐近线。在 \(x=1\) 处分母为 0 且不能约去，所以有垂直渐近线
\[
x=1.
\]
当 \(x\to\pm\infty\) 时，
\[
\frac{x^2+x}{x^2-1}\to1,
\]
故有水平渐近线
\[
y=1.
\]
共有 2 条。

\answer{C}
\examnote{分母为零不一定是垂直渐近线，能约去的点通常是可去间断点。}
\end{solutionblock}

\begin{problemblock}
\textbf{18.} 设曲线 \(y=f(x)\) 与 \(y=x^2-x\) 在点 \((1,0)\) 处有公共切线，则
\[
\lim_{n\to\infty}n f\left(\frac{n}{n+2}\right)=\underline{\qquad}.
\]
\end{problemblock}

\begin{solutionblock}
\analysis{公共切线给出 \(f(1)\) 与 \(f'(1)\)。}
曲线 \(y=x^2-x\) 在 \(x=1\) 处函数值为 0，导数为
\[
2x-1\big|_{x=1}=1.
\]
因为两曲线在 \((1,0)\) 处有公共切线，所以
\[
f(1)=0,\qquad f'(1)=1.
\]
令
\[
x_n=\frac{n}{n+2}=1-\frac{2}{n+2}.
\]
则
\[
x_n-1=-\frac{2}{n+2}.
\]
由可导性，
\[
f(x_n)=f(1)+f'(1)(x_n-1)+o(x_n-1)
=-\frac{2}{n+2}+o\left(\frac1n\right).
\]
因此
\[
n f\left(\frac{n}{n+2}\right)\to -2.
\]
\answer{\(-2\)}
\examnote{公共切线题常直接转化为函数值相同、导数相同。}
\end{solutionblock}

\section{填空题与计算题}

\begin{problemblock}
\textbf{19.} 已知
\[
f(x)=\frac{(x-1)(x-2)\cdots(x-n)}{(x+1)(x+2)\cdots(x+n)},
\]
则
\[
f'(1)=\underline{\qquad}.
\]
\end{problemblock}

\begin{solutionblock}
\analysis{因为 \(x=1\) 是分子的单根，求导值只需把 \((x-1)\) 提出来。}
写成
\[
f(x)=(x-1)\cdot \frac1{x+1}\prod_{k=2}^n\frac{x-k}{x+k}.
\]
因此
\[
f'(1)=\frac12\prod_{k=2}^n\frac{1-k}{1+k}.
\]
而
\[
\prod_{k=2}^n\frac{1-k}{1+k}
=(-1)^{n-1}\frac{1\cdot2\cdots(n-1)}{3\cdot4\cdots(n+1)}
=(-1)^{n-1}\frac{2}{n(n+1)}.
\]
故
\[
f'(1)=\frac{(-1)^{n-1}}{n(n+1)}.
\]
\answer{\(\displaystyle\frac{(-1)^{n-1}}{n(n+1)}\)}
\examnote{有简单零点时，\(f'(x_0)\) 等于去掉零因子后的剩余部分在 \(x_0\) 的值。}
\end{solutionblock}

\begin{problemblock}
\textbf{20.}（数学三不要求）曲线
\[
\begin{cases}
x=\displaystyle\int_0^{1-t}e^{-u^2}\,du,\\
y=t^2\ln(2-t^2)
\end{cases}
\]
在点 \((0,0)\) 处的切线方程为 \(\underline{\qquad}\)。
\end{problemblock}

\begin{solutionblock}
\analysis{先确定对应参数，再用 \(\frac{dy}{dx}=\frac{dy/dt}{dx/dt}\)。}
点 \((0,0)\) 对应
\[
x=0\Rightarrow 1-t=0\Rightarrow t=1,
\]
且此时
\[
y=1^2\ln1=0.
\]
求导：
\[
\frac{dx}{dt}=-e^{-(1-t)^2},
\]
故
\[
\left.\frac{dx}{dt}\right|_{t=1}=-1.
\]
又
\[
\frac{dy}{dt}=2t\ln(2-t^2)+t^2\frac{-2t}{2-t^2},
\]
所以
\[
\left.\frac{dy}{dt}\right|_{t=1}=-2.
\]
因此
\[
\left.\frac{dy}{dx}\right|_{t=1}=2.
\]
切线过 \((0,0)\)，方程为
\[
y=2x.
\]
\answer{\(y=2x\)}
\examnote{参数方程切线先找参数值，不要直接消参。}
\end{solutionblock}

\begin{problemblock}
\textbf{21.}（数学三不要求）对数螺线
\[
r=e^\theta
\]
在点
\[
(r,\theta)=\left(e^{\pi/2},\frac\pi2\right)
\]
处的切线的直角坐标方程为 \(\underline{\qquad}\)。
\end{problemblock}

\begin{solutionblock}
\analysis{极坐标曲线用 \(x=r\cos\theta,\ y=r\sin\theta\) 求斜率。}
对 \(r=e^\theta\)，有 \(r'=r\)。切线斜率公式为
\[
\frac{dy}{dx}
=\frac{r'\sin\theta+r\cos\theta}{r'\cos\theta-r\sin\theta}.
\]
在 \(\theta=\frac\pi2\) 处，
\[
\frac{dy}{dx}
=\frac{r}{-r}=-1.
\]
该点直角坐标为
\[
x=0,\qquad y=e^{\pi/2}.
\]
故切线方程为
\[
y-e^{\pi/2}=-x,
\]
即
\[
x+y=e^{\pi/2}.
\]
\answer{\(x+y=e^{\pi/2}\)}
\examnote{极坐标点要先转成直角坐标点，再写切线方程。}
\end{solutionblock}

\begin{problemblock}
\textbf{22.} 设函数
\[
f(x)=
\begin{cases}
\ln\sqrt{x},&x\ge1,\\
2x-1,&x<1,
\end{cases}
\qquad y=f(f(x)),
\]
则
\[
\left.\frac{dy}{dx}\right|_{x=e}
=\underline{\qquad}.
\]
\end{problemblock}

\begin{solutionblock}
\analysis{复合函数求导时要先判断内层函数值落在哪个分段。}
因为 \(e>1\)，
\[
f(e)=\ln\sqrt e=\frac12<1.
\]
所以外层函数在输入 \(\frac12\) 时使用分段
\[
f(u)=2u-1,\qquad f'(u)=2.
\]
内层在 \(x=e\) 处使用
\[
f(x)=\ln\sqrt{x}=\frac12\ln x,
\]
故
\[
f'(e)=\frac1{2e}.
\]
由链式法则
\[
\left.\frac{dy}{dx}\right|_{x=e}
=f'(f(e))f'(e)=2\cdot\frac1{2e}=\frac1e.
\]
\answer{\(\frac1e\)}
\examnote{分段复合函数一定要先算内层值，再决定外层用哪一段。}
\end{solutionblock}

\begin{problemblock}
\textbf{23.} 设 \(y=f(x)\) 的反函数是 \(x=\varphi(y)\)，且
\[
f(x)=\int_1^{2x}e^{t^2}\,dt+1,
\]
则
\[
\varphi''(1)=\underline{\qquad}.
\]
\end{problemblock}

\begin{solutionblock}
\analysis{反函数二阶导公式为 \(\varphi''(y)=-\frac{f''(x)}{[f'(x)]^3}\)，其中 \(x=\varphi(y)\)。}
先求 \(\varphi(1)\)。由
\[
f(x)=1
\]
得积分为 0，所以
\[
2x=1,\qquad x=\frac12.
\]
即
\[
\varphi(1)=\frac12.
\]
由变上限积分求导，
\[
f'(x)=2e^{(2x)^2}=2e^{4x^2},
\]
\[
f''(x)=16xe^{4x^2}.
\]
在 \(x=\frac12\) 处，
\[
f'\left(\frac12\right)=2e,\qquad
f''\left(\frac12\right)=8e.
\]
因此
\[
\varphi''(1)=
-\frac{8e}{(2e)^3}
=-\frac1{e^2}.
\]
\answer{\(-e^{-2}\)}
\examnote{反函数求导题先找对应的 \(x\)，再套公式。}
\end{solutionblock}

\begin{problemblock}
\textbf{24.} 函数
\[
y=x\ln(1-2x)
\]
在 \(x=0\) 处的 \(n(n\ge2)\) 阶导数
\[
y^{(n)}(0)=\underline{\qquad}.
\]
\end{problemblock}

\begin{solutionblock}
\analysis{用幂级数读出 \(x^n\) 的系数。}
当 \(|x|<\frac12\) 时，
\[
\ln(1-2x)=-\sum_{k=1}^{\infty}\frac{(2x)^k}{k}.
\]
因此
\[
y=x\ln(1-2x)
=-\sum_{k=1}^{\infty}\frac{2^k}{k}x^{k+1}.
\]
要取 \(x^n\) 项，令 \(k+1=n\)，即 \(k=n-1\)。故 \(x^n\) 的系数为
\[
-\frac{2^{n-1}}{n-1}.
\]
所以
\[
y^{(n)}(0)=n!\left(-\frac{2^{n-1}}{n-1}\right)
=-\frac{n!\,2^{n-1}}{n-1}.
\]
\answer{\(\displaystyle-\frac{n!\,2^{n-1}}{n-1}\)}
\examnote{高阶导数在 0 处等于 \(n!\) 乘 Taylor 展开中 \(x^n\) 的系数。}
\end{solutionblock}

\begin{problemblock}
\textbf{25.} 设
\[
f(x)=\frac{x^2+x-1}{x^2+x-2},
\]
则
\[
f^{(n)}(x)=\underline{\qquad}.
\]
\end{problemblock}

\begin{solutionblock}
\analysis{先把有理函数拆成简单分式。}
因为
\[
x^2+x-1=(x^2+x-2)+1,
\]
所以
\[
f(x)=1+\frac1{x^2+x-2}
=1+\frac1{(x-1)(x+2)}.
\]
部分分式分解：
\[
\frac1{(x-1)(x+2)}
=\frac13\left(\frac1{x-1}-\frac1{x+2}\right).
\]
当 \(n\ge1\) 时，常数项导数为 0，且
\[
\frac{d^n}{dx^n}\frac1{x-a}
=(-1)^n n!(x-a)^{-n-1}.
\]
故
\[
f^{(n)}(x)=\frac{(-1)^n n!}{3}
\left[
\frac1{(x-1)^{n+1}}-\frac1{(x+2)^{n+1}}
\right].
\]
\answer{\(\displaystyle \frac{(-1)^n n!}{3}\left[\frac1{(x-1)^{n+1}}-\frac1{(x+2)^{n+1}}\right]\)}
\examnote{有理函数高阶导优先拆成 \((x-a)^{-1}\) 的线性组合。}
\end{solutionblock}

\begin{problemblock}
\textbf{26.} 函数
\[
f(x)=\ln\bigl[(x-1)(x-2)\cdots(x-n)\bigr]
\]
的驻点个数为 \(\underline{\qquad}\)。
\end{problemblock}

\begin{solutionblock}
\analysis{驻点满足 \(f'(x)=0\)，并且只能在函数有定义的区间内讨论。}
求导得
\[
f'(x)=\sum_{k=1}^n\frac1{x-k}.
\]
再求导：
\[
f''(x)=-\sum_{k=1}^n\frac1{(x-k)^2}<0.
\]
所以 \(f'(x)\) 在每个定义区间上严格递减。

在区间 \((k,k+1)\ (k=1,2,\cdots,n-1)\) 上，
\[
\lim_{x\to k^+}f'(x)=+\infty,\qquad
\lim_{x\to(k+1)^-}f'(x)=-\infty.
\]
由连续性和严格递减性，每个这样的区间内恰有一个零点。

在 \((-\infty,1)\) 上，各项均为负，无零点；在 \((n,+\infty)\) 上，各项均为正，也无零点。因此驻点个数为
\[
n-1.
\]
\answer{\(n-1\)}
\examnote{对数乘积求导变成倒数和；再用单调性判断零点个数。}
\end{solutionblock}

\begin{problemblock}
\textbf{27.} 已知方程
\[
x^4+2x^3-3x^2-4x+a=0
\]
有两个重根，则
\[
a=\underline{\qquad}.
\]
\end{problemblock}

\begin{solutionblock}
\analysis{重根必为多项式与导函数的公共根。}
设
\[
P(x)=x^4+2x^3-3x^2-4x+a.
\]
则
\[
P'(x)=4x^3+6x^2-6x-4
=2(x-1)(2x+1)(x+2).
\]
因此重根只能在
\[
x=1,\quad x=-\frac12,\quad x=-2
\]
中产生。分别计算不含 \(a\) 的部分
\[
Q(x)=x^4+2x^3-3x^2-4x.
\]
有
\[
Q(1)=-4,\qquad Q(-2)=-4,
\]
所以若 \(a=4\)，则 \(x=1\) 与 \(x=-2\) 都是重根。事实上
\[
x^4+2x^3-3x^2-4x+4=(x-1)^2(x+2)^2.
\]
故
\[
a=4.
\]
\answer{4}
\examnote{“两个重根”对四次方程通常意味着可以分解成两个平方因子。}
\end{solutionblock}

\begin{problemblock}
\textbf{28.} 已知方程
\[
3x^4-8x^3-6x^2+24x+a=0
\]
有四个不同的实根，则 \(a\) 的取值范围为 \(\underline{\qquad}\)。
\end{problemblock}

\begin{solutionblock}
\analysis{把参数 \(a\) 看作竖直平移，先研究不含 \(a\) 的四次函数极值。}
令
\[
g(x)=3x^4-8x^3-6x^2+24x.
\]
则
\[
g'(x)=12x^3-24x^2-12x+24
=12(x+1)(x-1)(x-2).
\]
三个驻点为
\[
x=-1,\quad x=1,\quad x=2.
\]
计算函数值：
\[
g(-1)=-19,\qquad g(1)=13,\qquad g(2)=8.
\]
因为最高次系数为正，曲线从左到右依次为：下降到 \(-19\)，上升到 \(13\)，下降到 \(8\)，再上升到 \(+\infty\)。

方程
\[
g(x)+a=0
\]
有四个不同实根，等价于水平线 \(y=-a\) 同时穿过四个单调区间，必须满足
\[
8<-a<13.
\]
因此
\[
-13<a<-8.
\]
\answer{\((-13,-8)\)}
\examnote{四次方程实根个数常转化为函数图像与水平线交点个数。}
\end{solutionblock}

\begin{problemblock}
\textbf{29.} 设 \(f(x)\) 为连续函数，
\[
\lim_{x\to0}\frac{xf(x)-\ln(1+x)}{x^2}=2,
\]
\[
F(x)=\int_0^x t f(x-t)\,dt.
\]
当 \(x\to0\) 时，\(F(x)-\frac12x^2\) 与 \(bx^k\) 为等价无穷小，其中常数 \(b\ne0\)，\(k\) 为某正整数。求 \(k,b\) 及 \(f(0),f'(0)\)。
\end{problemblock}

\begin{solutionblock}
\analysis{先由已知极限反求 \(f(x)\) 的一阶展开，再代入积分。}
由
\[
\ln(1+x)=x-\frac{x^2}{2}+o(x^2),
\]
可得
\[
xf(x)=\ln(1+x)+2x^2+o(x^2)
=x+\frac32x^2+o(x^2).
\]
因此
\[
f(x)=1+\frac32x+o(x),
\]
所以
\[
f(0)=1,\qquad f'(0)=\frac32.
\]
再看
\[
F(x)=\int_0^x t f(x-t)\,dt.
\]
当 \(x\to0\) 时，
\[
f(x-t)=1+\frac32(x-t)+o(x)
\]
在积分区间内成立。于是
\[
F(x)=\int_0^x t\,dt
+\frac32\int_0^x t(x-t)\,dt
+o(x^3).
\]
计算
\[
\int_0^x t\,dt=\frac{x^2}{2},
\]
\[
\int_0^x t(x-t)\,dt
=x\cdot\frac{x^2}{2}-\frac{x^3}{3}
=\frac{x^3}{6}.
\]
因此
\[
F(x)=\frac{x^2}{2}+\frac32\cdot\frac{x^3}{6}+o(x^3)
=\frac{x^2}{2}+\frac{x^3}{4}+o(x^3).
\]
故
\[
F(x)-\frac12x^2\sim\frac14x^3.
\]
\answer{\(k=3,\ b=\frac14,\ f(0)=1,\ f'(0)=\frac32\)}
\examnote{积分里 \(f(x-t)\) 的一阶项会贡献三阶小量，这是本题的主部来源。}
\end{solutionblock}

\begin{problemblock}
\textbf{30.} 已知函数 \(f(u)\) 具有二阶导数，且 \(f'(0)=1\)，函数 \(y=y(x)\) 由方程
\[
y-xe^{y-1}=1
\]
所确定。设
\[
z=f(\ln y-\sin x),
\]
求
\[
\left.\frac{dz}{dx}\right|_{x=0},\qquad
\left.\frac{d^2z}{dx^2}\right|_{x=0}.
\]
\end{problemblock}

\begin{solutionblock}
\analysis{先求隐函数 \(y\) 在 0 处的一、二阶导数，再处理复合函数。}
由方程在 \(x=0\) 处得
\[
y(0)=1.
\]
设
\[
y=1+ax+bx^2+o(x^2).
\]
则
\[
y-1=ax+bx^2+o(x^2),
\]
\[
e^{y-1}=1+ax+\left(b+\frac{a^2}{2}\right)x^2+o(x^2).
\]
代入
\[
y-xe^{y-1}=1
\]
得
\[
1+ax+bx^2-x(1+ax+o(x))=1.
\]
比较系数：
\[
a=1,\qquad b=a=1.
\]
所以
\[
y'(0)=1,\qquad y''(0)=2.
\]
令
\[
u(x)=\ln y-\sin x.
\]
则
\[
u(0)=0,
\]
\[
u'(0)=\frac{y'(0)}{y(0)}-\cos0=1-1=0.
\]
又
\[
u''(0)=\left.\left(\frac{y''y-(y')^2}{y^2}+\sin x\right)\right|_{x=0}
=2-1=1.
\]
因为
\[
z=f(u),
\]
所以
\[
z'(0)=f'(0)u'(0)=0,
\]
\[
z''(0)=f''(0)[u'(0)]^2+f'(0)u''(0)=1.
\]
\answer{\(\left.z'\right|_{0}=0,\quad \left.z''\right|_{0}=1\)}
\examnote{若内层函数 \(u'(0)=0\)，二阶复合求导中的 \(f''(0)[u'(0)]^2\) 项会消失。}
\end{solutionblock}

\begin{problemblock}
\textbf{31.}（数学三不要求）设 \(f(t)\) 二阶可导，且 \(f''(t)\ne0\)，
\[
\begin{cases}
y=tf'(t)-f(t),\\
x=f'(t),
\end{cases}
\]
求
\[
\frac{d^2y}{dx^2}\quad\text{及}\quad \frac{d^2x}{dy^2}.
\]
\end{problemblock}

\begin{solutionblock}
\analysis{参数方程求二阶导，先求 \(\frac{dy}{dx}\)，再对参数继续求导。}
由
\[
\frac{dx}{dt}=f''(t),
\]
\[
\frac{dy}{dt}=f'(t)+tf''(t)-f'(t)=tf''(t),
\]
得
\[
\frac{dy}{dx}=\frac{dy/dt}{dx/dt}=t.
\]
因此
\[
\frac{d^2y}{dx^2}
=\frac{d}{dx}\left(\frac{dy}{dx}\right)
=\frac{dt}{dx}
=\frac1{f''(t)}.
\]
又
\[
\frac{dx}{dy}=\frac1{dy/dx}=\frac1t.
\]
于是
\[
\frac{d^2x}{dy^2}
=\frac{d}{dy}\left(\frac1t\right)
=\frac{\frac{d}{dt}(1/t)}{dy/dt}
=\frac{-1/t^2}{tf''(t)}
=-\frac1{t^3f''(t)}.
\]
\answer{\(\displaystyle \frac{d^2y}{dx^2}=\frac1{f''(t)},\quad \frac{d^2x}{dy^2}=-\frac1{t^3f''(t)}\)}
\examnote{二阶参数导数本质是“先对 \(t\) 求导，再除以对应的一阶导”。}
\end{solutionblock}

\begin{problemblock}
\textbf{32.}（数学三不要求）设 \(y=y(x)\) 由
\[
\begin{cases}
x=t^3+2t+1,\\
t-\displaystyle\int_1^{y+t^2}e^{-u^2}\,du=0
\end{cases}
\]
确定，求
\[
\left.\frac{dy}{dx}\right|_{t=0},\qquad
\left.\frac{d^2y}{dx^2}\right|_{t=0}.
\]
\end{problemblock}

\begin{solutionblock}
\analysis{先把 \(y\) 看成 \(t\) 的函数，再用参数方程求导。}
当 \(t=0\) 时，第二个方程给出
\[
\int_1^y e^{-u^2}\,du=0,
\]
故
\[
y(0)=1.
\]
对
\[
t-\int_1^{y+t^2}e^{-u^2}\,du=0
\]
求导：
\[
1-e^{-(y+t^2)^2}(y'+2t)=0,
\]
即
\[
y'+2t=e^{(y+t^2)^2}.
\]
在 \(t=0\) 处，
\[
y'(0)=e.
\]
又
\[
x'=3t^2+2,\qquad x'(0)=2.
\]
所以
\[
\left.\frac{dy}{dx}\right|_{t=0}=\frac{y'(0)}{x'(0)}=\frac e2.
\]
继续对
\[
y'+2t=e^{(y+t^2)^2}
\]
求导。令 \(s=y+t^2\)，则右侧导数为
\[
e^{s^2}\cdot2s(y'+2t).
\]
在 \(t=0\) 处，\(s=1,\ y'+2t=e\)，故
\[
y''(0)+2=2e^2,
\]
\[
y''(0)=2e^2-2.
\]
而 \(x''(0)=0\)，故
\[
\left.\frac{d^2y}{dx^2}\right|_{t=0}
=\left.\frac{y''x'-y'x''}{(x')^3}\right|_{t=0}
=\frac{(2e^2-2)\cdot2}{8}
=\frac{e^2-1}{2}.
\]
\answer{\(\displaystyle \left.\frac{dy}{dx}\right|_{t=0}=\frac e2,\quad \left.\frac{d^2y}{dx^2}\right|_{t=0}=\frac{e^2-1}{2}\)}
\examnote{隐式参数题要分清：先求 \(dy/dt\)、\(d^2y/dt^2\)，再换成对 \(x\) 的导数。}
\end{solutionblock}

\begin{problemblock}
\textbf{33.} 设函数
\[
\varphi(x)=\int_0^{\sin x}f(tx^2)\,dt,
\]
其中 \(f(x)\) 是连续函数，且 \(f(0)=2\)。

(1) 求 \(\varphi'(x)\)；

(2) 讨论 \(\varphi'(x)\) 的连续性。
\end{problemblock}

\begin{solutionblock}
\analysis{由于只假设 \(f\) 连续，不宜直接对 \(f(tx^2)\) 求偏导。可先作变量代换。}
当 \(x\ne0\) 时，令 \(u=tx^2\)，得
\[
\varphi(x)=\frac1{x^2}\int_0^{x^2\sin x}f(u)\,du.
\]
记
\[
H(u)=\int_0^u f(s)\,ds,\qquad u=x^2\sin x.
\]
则
\[
\varphi(x)=x^{-2}H(x^2\sin x).
\]
所以当 \(x\ne0\) 时，
\[
\varphi'(x)
=-\frac2{x^3}H(x^2\sin x)
+\frac1{x^2}f(x^2\sin x)(2x\sin x+x^2\cos x).
\]
即
\[
\varphi'(x)=
-\frac{2}{x^3}\int_0^{x^2\sin x}f(u)\,du
+f(x^2\sin x)\left(\frac{2\sin x}{x}+\cos x\right).
\]
当 \(x=0\) 时，
\[
\varphi(x)=\int_0^{\sin x}f(tx^2)\,dt
\sim\int_0^{x}2\,dt=2x,
\]
故
\[
\varphi'(0)=2.
\]

连续性方面，当 \(x\to0\) 时，由 \(f\) 在 0 连续，\(H(u)\sim2u\)，且 \(x^2\sin x\sim x^3\)。代入上式可得
\[
\varphi'(x)\to -\frac2{x^3}\cdot 2x^3+2(2+1)=2.
\]
因此 \(\varphi'(x)\) 在 \(x=0\) 处连续；在 \(x\ne0\) 处由连续函数复合可知 \(\varphi'(x)\) 连续。故 \(\varphi'(x)\) 在其定义域内连续。
\answer{见上式，且 \(\varphi'(x)\) 连续。}
\examnote{只有连续性时，优先用变上限积分函数 \(H\) 规避对 \(f\) 求导。}
\end{solutionblock}

\begin{problemblock}
\textbf{34.} 设 \(f(x)\) 连续，
\[
\varphi(x)=\int_0^1 f(xt)\,dt,
\]
且
\[
\lim_{x\to0}\frac{f(x)}x=A\quad(A\text{ 为常数}).
\]
求 \(\varphi'(x)\)，并讨论 \(\varphi'(x)\) 在 \(x=0\) 处的连续性。
\end{problemblock}

\begin{solutionblock}
\analysis{同样用变量代换，把含参积分化为变上限积分。}
由已知极限可知 \(f(0)=0\)。当 \(x\ne0\) 时，令 \(u=xt\)，
\[
\varphi(x)=\frac1x\int_0^x f(u)\,du.
\]
因此
\[
\varphi'(x)=\frac{x f(x)-\int_0^x f(u)\,du}{x^2}\qquad(x\ne0).
\]
当 \(x=0\) 时，
\[
\varphi'(0)=\lim_{x\to0}\frac{\varphi(x)-\varphi(0)}x
=\lim_{x\to0}\frac1{x^2}\int_0^x f(u)\,du.
\]
因为
\[
f(u)\sim Au,
\]
所以
\[
\int_0^x f(u)\,du\sim\frac A2x^2.
\]
故
\[
\varphi'(0)=\frac A2.
\]
再看连续性：
\[
x f(x)\sim Ax^2,\qquad
\int_0^x f(u)\,du\sim\frac A2x^2.
\]
于是
\[
\lim_{x\to0}\varphi'(x)
=\frac A2
=\varphi'(0).
\]
所以 \(\varphi'(x)\) 在 \(x=0\) 处连续。
\answer{\(\displaystyle \varphi'(x)=\frac{x f(x)-\int_0^x f(u)\,du}{x^2}(x\ne0),\ \varphi'(0)=\frac A2\)，且在 0 处连续。}
\examnote{\(\int_0^1f(xt)dt\) 是常见模型，换元后就是平均值函数。}
\end{solutionblock}

\begin{problemblock}
\textbf{35.} 设函数由方程
\[
2y^3-2y^2+2xy-x^2=1
\]
所确定，试求 \(y=y(x)\) 的驻点，并判别它是否为极值点。
\end{problemblock}

\begin{solutionblock}
\analysis{隐函数驻点满足 \(y'=0\)。}
令
\[
F(x,y)=2y^3-2y^2+2xy-x^2-1=0.
\]
则
\[
F_x=2y-2x,\qquad
F_y=6y^2-4y+2x.
\]
由隐函数求导，
\[
y'=-\frac{F_x}{F_y}
=\frac{2x-2y}{6y^2-4y+2x}.
\]
驻点要求 \(y'=0\)，即
\[
x=y.
\]
代入原方程：
\[
2y^3-2y^2+2y^2-y^2=1,
\]
即
\[
2y^3-y^2=1.
\]
整理得
\[
2y^3-y^2-1=(y-1)(2y^2+y+1)=0.
\]
二次因子无实根，所以
\[
y=1,\qquad x=1.
\]
驻点为 \((1,1)\)。

判别极值：在 \((1,1)\) 处，
\[
F_y=4\ne0,\qquad y'=0.
\]
二阶求导公式在驻点处化为
\[
F_{xx}+F_y y''=0.
\]
由于 \(F_{xx}=-2\)，故
\[
-2+4y''=0,\qquad y''=\frac12>0.
\]
因此 \((1,1)\) 是极小值点。
\answer{驻点为 \((1,1)\)，且为极小值点。}
\examnote{隐函数驻点先令 \(F_x=0\)，再代回原方程；判别时用二阶导。}
\end{solutionblock}

\begin{problemblock}
\textbf{36.}（数学三不要求）已知曲线 \(L\) 的方程为
\[
\begin{cases}
x=t^2+1,\\
y=4t-t^2
\end{cases}
\quad(t\ge0).
\]
(1) 讨论 \(L\) 的凹凸性；

(2) 过点 \((-1,0)\) 引 \(L\) 的切线，求切点 \((x_0,y_0)\)，并写出切线方程；

(3) 求此切线与 \(L\)（对应于 \(x\le x_0\) 的部分）及 \(x\) 轴所围成的平面图形的面积。
\end{problemblock}

\begin{solutionblock}
\analysis{先用参数方程求斜率和凹凸，再由“切线过定点”确定参数。}
(1) 有
\[
\frac{dx}{dt}=2t,\qquad \frac{dy}{dt}=4-2t.
\]
当 \(t>0\) 时，
\[
\frac{dy}{dx}=\frac{4-2t}{2t}=\frac2t-1.
\]
所以
\[
\frac{d^2y}{dx^2}
=\frac{d}{dt}\left(\frac2t-1\right)\bigg/\frac{dx}{dt}
=\frac{-2/t^2}{2t}
=-\frac1{t^3}<0.
\]
故曲线凹向下。

(2) 参数为 \(t\) 的点为
\[
(x,y)=(t^2+1,\,4t-t^2),
\]
切线斜率为
\[
m=\frac2t-1.
\]
若切线过 \((-1,0)\)，则
\[
\frac{4t-t^2}{t^2+1-(-1)}
=\frac{4t-t^2}{t^2+2}
=\frac2t-1.
\]
解得
\[
t=1\quad(t\ge0).
\]
故切点
\[
(x_0,y_0)=(2,3),
\]
斜率
\[
m=1.
\]
切线方程为
\[
y=x+1.
\]

(3) 所求面积等于直线 \(y=x+1\) 在 \([-1,2]\) 上方与 \(x\) 轴围成的三角形面积，减去曲线 \(L\) 在 \(t\in[0,1]\) 下方的面积：
\[
S=\int_{-1}^{2}(x+1)\,dx-\int_{0}^{1}y(t)x'(t)\,dt.
\]
第一项
\[
\int_{-1}^{2}(x+1)\,dx=\frac92.
\]
第二项
\[
\int_0^1(4t-t^2)\cdot2t\,dt
=\int_0^1(8t^2-2t^3)\,dt
=\frac83-\frac12=\frac{13}{6}.
\]
因此
\[
S=\frac92-\frac{13}{6}
=\frac{7}{3}.
\]
\answer{(1) 凹向下；(2) 切点 \((2,3)\)，切线 \(y=x+1\)；(3) 面积 \(\frac73\)。}
\examnote{参数曲线面积常用 \(\int y\,dx=\int y(t)x'(t)\,dt\)。}
\end{solutionblock}

\begin{problemblock}
\textbf{37.} 试确定方程
\[
x^3-x=\sin x
\]
的实根个数。
\end{problemblock}

\begin{solutionblock}
\analysis{方程两边移项，利用奇偶性和导数研究正半轴。}
令
\[
F(x)=x^3-x-\sin x.
\]
则 \(F(x)\) 是奇函数，且
\[
F(0)=0.
\]
只需研究 \(x>0\)。有
\[
F'(x)=3x^2-1-\cos x,
\]
\[
F''(x)=6x+\sin x>0\qquad(x>0).
\]
所以 \(F'(x)\) 在 \((0,+\infty)\) 上严格递增。又
\[
F'(0)=-2,\qquad F'(x)\to+\infty,
\]
故 \(F'\) 在正半轴只有一个零点，\(F\) 先减后增。

当 \(x\to0^+\) 时，\(F(x)<0\)；当 \(x\to+\infty\) 时，\(F(x)\to+\infty\)。因此 \(F(x)=0\) 在 \((0,+\infty)\) 上恰有一个根。由奇函数对称性，在 \((-\infty,0)\) 上也恰有一个根，再加上 \(x=0\)，共有 3 个实根。
\answer{3 个}
\examnote{奇函数方程常先找 \(0\) 根，再研究正半轴。}
\end{solutionblock}

\begin{problemblock}
\textbf{38.} 试确定方程
\[
\int_0^x e^{-t^2}\,dt=x^3-x
\]
的实根个数。
\end{problemblock}

\begin{solutionblock}
\analysis{移项后仍是奇函数，研究正半轴即可。}
令
\[
F(x)=\int_0^x e^{-t^2}\,dt-x^3+x.
\]
则 \(F\) 为奇函数，且 \(F(0)=0\)。对 \(x>0\)，
\[
F'(x)=e^{-x^2}-3x^2+1,
\]
\[
F''(x)=-2xe^{-x^2}-6x<0.
\]
所以 \(F'(x)\) 在 \((0,+\infty)\) 上严格递减。又
\[
F'(0)=2,\qquad F'(x)\to-\infty,
\]
故 \(F'\) 在正半轴仅有一个零点，\(F\) 先增后减。

当 \(x\to0^+\) 时，\(F(x)>0\)；当 \(x\to+\infty\) 时，\(F(x)\to-\infty\)。所以正半轴上恰有一个零点。由奇函数对称性，负半轴上也恰有一个零点，再加 \(x=0\)，共有 3 个实根。
\answer{3 个}
\examnote{“导数单调”能控制函数最多穿过几次横轴。}
\end{solutionblock}

\begin{problemblock}
\textbf{39.} 试确定方程
\[
e^x=ax^2\quad(a>0)
\]
的实根个数。
\end{problemblock}

\begin{solutionblock}
\analysis{把参数写成函数值：\(a=\frac{e^x}{x^2}\)，注意 \(x=0\) 不可能是根。}
令
\[
\phi(x)=\frac{e^x}{x^2}\qquad(x\ne0).
\]
则方程等价于
\[
a=\phi(x).
\]
在 \(x<0\) 上，
\[
\frac{\phi'(x)}{\phi(x)}=1-\frac2x>0,
\]
所以 \(\phi\) 严格递增，且
\[
\phi(x)\to0\ (x\to-\infty),\qquad \phi(x)\to+\infty\ (x\to0^-).
\]
故负半轴上对任意 \(a>0\) 都有一个根。

在 \(x>0\) 上，\(\phi\) 在 \(x=2\) 处取得最小值
\[
\phi(2)=\frac{e^2}{4}.
\]
因此正半轴上：
若 \(0<a<\frac{e^2}{4}\)，无根；若 \(a=\frac{e^2}{4}\)，有一个根；若 \(a>\frac{e^2}{4}\)，有两个根。

综上：
\[
\begin{cases}
0<a<\frac{e^2}{4},&1\text{ 个实根},\\
a=\frac{e^2}{4},&2\text{ 个实根},\\
a>\frac{e^2}{4},&3\text{ 个实根}.
\end{cases}
\]
\answer{见上分段结论}
\examnote{含参数根个数题，常转化为参数等于某函数值，再研究该函数值域。}
\end{solutionblock}

\begin{problemblock}
\textbf{40.} 试确定方程
\[
\ln x=kx
\]
的实根个数。
\end{problemblock}

\begin{solutionblock}
\analysis{定义域为 \(x>0\)，化为 \(\frac{\ln x}{x}=k\)。}
令
\[
\phi(x)=\frac{\ln x}{x}\quad(x>0).
\]
则
\[
\phi'(x)=\frac{1-\ln x}{x^2}.
\]
所以 \(\phi\) 在 \((0,e)\) 上递增，在 \((e,+\infty)\) 上递减，最大值为
\[
\phi(e)=\frac1e.
\]
并且
\[
\phi(x)\to-\infty\quad(x\to0^+),\qquad
\phi(x)\to0^+\quad(x\to+\infty).
\]
因此：
\[
\begin{cases}
k>\frac1e,&0\text{ 个实根},\\
k=\frac1e,&1\text{ 个实根},\\
0<k<\frac1e,&2\text{ 个实根},\\
k=0,&1\text{ 个实根},\\
k<0,&1\text{ 个实根}.
\end{cases}
\]
\answer{见上分段结论}
\examnote{\(\ln x=kx\) 是经典模型，关键函数是 \(\ln x/x\)。}
\end{solutionblock}

\begin{problemblock}
\textbf{41.} 试证：当 \(x\ge0\) 时，
\[
x\le e^x\ln(1+x).
\]
\end{problemblock}

\begin{solutionblock}
\analysis{构造差函数并证明单调递增。}
令
\[
F(x)=e^x\ln(1+x)-x.
\]
则
\[
F(0)=0.
\]
求导：
\[
F'(x)=e^x\ln(1+x)+\frac{e^x}{1+x}-1.
\]
当 \(x\ge0\) 时，
\[
e^x\ge1+x,
\]
故
\[
\frac{e^x}{1+x}\ge1.
\]
又 \(e^x\ln(1+x)\ge0\)，所以
\[
F'(x)\ge0.
\]
因此 \(F(x)\ge F(0)=0\)，即
\[
x\le e^x\ln(1+x).
\]
\examnote{证明不等式时，若右边含指数与对数，构造差函数后常用基本不等式 \(e^x\ge1+x\)。}
\end{solutionblock}

\begin{problemblock}
\textbf{42.} 设 \(x>0\)，证明：
\[
2\sin x+e^x-e^{-x}>4x.
\]
\end{problemblock}

\begin{solutionblock}
\analysis{把不等式化为函数正性，用二阶导控制单调性。}
令
\[
F(x)=\sin x+\sinh x-2x.
\]
原不等式等价于
\[
2F(x)>0.
\]
有
\[
F(0)=0,\qquad F'(0)=\cos0+\cosh0-2=0.
\]
再求二阶导：
\[
F''(x)=-\sin x+\sinh x.
\]
当 \(x>0\) 时，\(\sinh x>x\)，且 \(\sin x\le x\)，所以
\[
F''(x)>0.
\]
于是 \(F'(x)>F'(0)=0\)，进而
\[
F(x)>F(0)=0.
\]
故
\[
2\sin x+e^x-e^{-x}>4x.
\]
\examnote{若一阶导在 0 处也为 0，可继续看二阶导的符号。}
\end{solutionblock}

\begin{problemblock}
\textbf{43.} 设 \(x>0\)，常数 \(a>e\)。证明：
\[
(a+x)^a<a^{a+x}.
\]
\end{problemblock}

\begin{solutionblock}
\analysis{两边取对数，将幂不等式化为对数不等式。}
原不等式等价于
\[
a\ln(a+x)<(a+x)\ln a.
\]
两边除以 \(a>0\)，令 \(t=\frac{x}{a}>0\)，得
\[
\ln a+\ln(1+t)<(1+t)\ln a,
\]
即
\[
\ln(1+t)<t\ln a.
\]
由于
\[
\ln(1+t)<t\qquad(t>0),
\]
且 \(a>e\Rightarrow \ln a>1\)，所以
\[
\ln(1+t)<t<t\ln a.
\]
结论成立。
\examnote{幂不等式优先取对数；出现 \(a+x\) 时，令 \(t=x/a\) 可简化结构。}
\end{solutionblock}

\begin{problemblock}
\textbf{44.} 设 \(e<a<b\)，证明：
\[
a^2<ab\frac{\ln a}{\ln b}<b^2.
\]
\end{problemblock}

\begin{solutionblock}
\analysis{拆成左右两个不等式，分别转化为单调函数。}
左边不等式
\[
a^2<ab\frac{\ln a}{\ln b}
\]
等价于
\[
a\ln b<b\ln a,
\]
即
\[
\frac{\ln b}{b}<\frac{\ln a}{a}.
\]
函数
\[
\phi(x)=\frac{\ln x}{x}
\]
在 \(x>e\) 上单调递减。由于 \(e<a<b\)，故
\[
\frac{\ln a}{a}>\frac{\ln b}{b},
\]
左边不等式成立。

右边不等式
\[
ab\frac{\ln a}{\ln b}<b^2
\]
等价于
\[
a\ln a<b\ln b.
\]
函数
\[
\psi(x)=x\ln x
\]
在 \(x>e\) 上单调递增，因此
\[
a\ln a<b\ln b.
\]
右边不等式成立。故原不等式成立。
\examnote{同一个双边不等式两侧可能对应不同的单调函数。}
\end{solutionblock}

\begin{problemblock}
\textbf{45.} 设 \(f(x)\) 和 \(g(x)\) 在 \([0,1]\) 上连续，在 \((0,1)\) 内可导，
\[
f(0)=f(1)=-1,\qquad \int_0^1 f(x)\,dx>\frac12.
\]
试证至少存在一点 \(\xi\in(0,1)\)，使
\[
f'(\xi)+g'(\xi)[f(\xi)-\xi]=1.
\]
\end{problemblock}

\begin{solutionblock}
\analysis{目标式可看成某个辅助函数导数为 0。注意
\[
\frac{d}{dx}\left(e^{g(x)}(f(x)-x)\right)
=e^{g(x)}\{f'(x)-1+g'(x)[f(x)-x]\}.
\]
}
令
\[
H(x)=e^{g(x)}(f(x)-x).
\]
若能证明 \(H\) 在 \((0,1)\) 内有极值点，则由费马定理 \(H'(\xi)=0\)，即
\[
f'(\xi)+g'(\xi)[f(\xi)-\xi]=1.
\]
下面证明内点极值存在。由
\[
\int_0^1 f(x)\,dx>\frac12
\]
得
\[
\int_0^1 [f(x)-x]\,dx>0.
\]
因此存在 \(c\in(0,1)\)，使
\[
f(c)-c>0.
\]
于是
\[
H(c)>0.
\]
而
\[
H(0)=e^{g(0)}(f(0)-0)=-e^{g(0)}<0,
\]
\[
H(1)=e^{g(1)}(f(1)-1)=-2e^{g(1)}<0.
\]
所以连续函数 \(H\) 在 \([0,1]\) 上的最大值必在某个内点 \(\xi\in(0,1)\) 取得。由费马定理，
\[
H'(\xi)=0.
\]
即
\[
e^{g(\xi)}\{f'(\xi)-1+g'(\xi)[f(\xi)-\xi]\}=0.
\]
由于 \(e^{g(\xi)}>0\)，故
\[
f'(\xi)+g'(\xi)[f(\xi)-\xi]=1.
\]
\examnote{看到 \(f'+g'(f-x)\) 时，尝试构造 \(e^g(f-x)\)；这是“配导数”的典型技巧。}
\end{solutionblock}

\begin{problemblock}
\textbf{46.} 设 \(f(x),g(x)\) 在 \([0,1]\) 上连续，在 \((0,1)\) 内可导，且
\[
\int_0^1 f(x)\,dx=3\int_{2/3}^1 f(x)\,dx.
\]
试证存在 \(\xi,\eta\in(0,1)\)，使
\[
f'(\xi)=g'(\xi)[f(\eta)-f(\xi)].
\]
\end{problemblock}

\begin{solutionblock}
\analysis{题干中 \(g\) 没有额外条件，因此关键不是研究 \(g\)，而是先由积分条件推出存在一点 \(f'(\xi)=0\)。原题没有要求 \(\xi,\eta\) 不同，于是可取 \(\eta=\xi\)。}
由题设
\[
\int_0^1 f(x)\,dx
=\int_0^{2/3}f(x)\,dx+\int_{2/3}^1f(x)\,dx
=3\int_{2/3}^1f(x)\,dx,
\]
得
\[
\int_0^{2/3}f(x)\,dx=2\int_{2/3}^1f(x)\,dx.
\]
于是
\[
\frac{1}{2/3}\int_0^{2/3}f(x)\,dx
=3\int_{2/3}^1f(x)\,dx
=\frac{1}{1/3}\int_{2/3}^1f(x)\,dx.
\]
由积分中值定理，存在
\[
u\in\left(0,\frac23\right),\qquad v\in\left(\frac23,1\right),
\]
使
\[
f(u)=\frac{1}{2/3}\int_0^{2/3}f(x)\,dx,\qquad
f(v)=\frac{1}{1/3}\int_{2/3}^1f(x)\,dx.
\]
因此
\[
f(u)=f(v).
\]
函数 \(f\) 在 \([u,v]\) 上连续、在 \((u,v)\) 内可导，由罗尔定理，存在
\[
\xi\in(u,v)\subset(0,1)
\]
使
\[
f'(\xi)=0.
\]
取
\[
\eta=\xi.
\]
则
\[
g'(\xi)[f(\eta)-f(\xi)]
=g'(\xi)[f(\xi)-f(\xi)]
=0=f'(\xi).
\]
故结论成立。
\examnote{遇到“只有 \(f\) 的条件却含 \(g'\)”时，要检查题目是否允许取 \(\eta=\xi\)。若允许，先造出 \(f'=0\) 往往即可。}
\end{solutionblock}

\begin{problemblock}
\textbf{47.} 设 \(f(x)\) 在 \([-2,2]\) 上二阶可导，且 \(|f(x)|\le 1\)，又
\[
[f(0)]^2+[f'(0)]^2=4.
\]
证明在 \((-2,2)\) 内至少存在一点 \(\xi\)，使
\[
f''(\xi)+f(\xi)=0.
\]
\end{problemblock}

\begin{solutionblock}
\analysis{反证法。把 \(f''+f\) 看成“受迫项”，构造
\[
F(x)=f'(x)\cos x+f(x)\sin x,
\]
则 \(F'(x)=[f''(x)+f(x)]\cos x\)。在 \((-\pi/2,\pi/2)\) 内 \(\cos x>0\)，可得到单调性。}
由于 \(|f(0)|\le 1\)，而
\[
[f(0)]^2+[f'(0)]^2=4,
\]
故
\[
|f'(0)|=\sqrt{4-[f(0)]^2}\ge \sqrt3>1.
\]
反设对任意 \(x\in(-2,2)\)，均有
\[
f''(x)+f(x)\ne 0.
\]
导函数具有介值性，所以 \(f''+f\) 在区间内不能由正跳到负，否则中间会取零。因此 \(f''+f\) 在 \((-2,2)\) 内同号。

令
\[
F(x)=f'(x)\cos x+f(x)\sin x.
\]
则
\[
F'(x)=[f''(x)+f(x)]\cos x.
\]
注意
\[
\left[-\frac{\pi}{2},\frac{\pi}{2}\right]\subset(-2,2),\qquad \cos x>0\quad\left(|x|<\frac{\pi}{2}\right).
\]
若 \(f''+f>0\)，则 \(F\) 在 \((-\pi/2,\pi/2)\) 上严格递增。
当 \(f'(0)>1\) 时，
\[
F\left(\frac{\pi}{2}\right)>F(0)=f'(0)>1,
\]
但
\[
F\left(\frac{\pi}{2}\right)=f\left(\frac{\pi}{2}\right),
\]
与 \(|f|\le1\) 矛盾。
当 \(f'(0)<-1\) 时，
\[
F\left(-\frac{\pi}{2}\right)<F(0)=f'(0)<-1,
\]
而
\[
F\left(-\frac{\pi}{2}\right)=-f\left(-\frac{\pi}{2}\right),
\]
也与 \(|f|\le1\) 矛盾。

若 \(f''+f<0\)，则 \(F\) 在 \((-\pi/2,\pi/2)\) 上严格递减。类似地：
当 \(f'(0)>1\) 时，
\[
F\left(-\frac{\pi}{2}\right)>F(0)=f'(0)>1,
\]
与 \(|f|\le1\) 矛盾；当 \(f'(0)<-1\) 时，
\[
F\left(\frac{\pi}{2}\right)<F(0)=f'(0)<-1,
\]
也与 \(|f|\le1\) 矛盾。

所有情形均矛盾，故必存在 \(\xi\in(-2,2)\)，使
\[
f''(\xi)+f(\xi)=0.
\]
\examnote{本题核心是把二阶表达式 \(f''+f\) 配成一阶函数的导数；\(\sin,\cos\) 辅助函数是处理 \(f''+f\) 的常用模板。}
\end{solutionblock}

\begin{problemblock}
\textbf{48.} 设函数 \(f(x)\) 在闭区间 \([a,b]\) 上连续，在开区间 \((a,b)\) 内可导，且 \(f'(x)>0\)。若极限
\[
\lim_{x\to a^+}\frac{f(2x-a)}{x-a}
\]
存在，证明：
\begin{enumerate}
\item 在 \((a,b)\) 内 \(f(x)>0\)；
\item 在 \((a,b)\) 内存在点 \(\xi\)，使
\[
\frac{b^2-a^2}{\int_a^b f(x)\,dx}=\frac{2\xi}{f(\xi)};
\]
\item 在 \((a,b)\) 内存在与 \((2)\) 中 \(\xi\) 相异的点 \(\eta\)，使
\[
f'(\eta)(b^2-a^2)=\frac{2\xi}{\xi-a}\int_a^b f(x)\,dx.
\]
\end{enumerate}
\end{problemblock}

\begin{solutionblock}
\analysis{先由极限存在推出 \(f(a)=0\)，再用 \(f'>0\) 得正性。后两问分别是柯西中值定理和拉格朗日中值定理。}
\textbf{(1)} 因为 \(f\) 在 \(a\) 处右连续，故
\[
\lim_{x\to a^+}f(2x-a)=f(a).
\]
若 \(f(a)\ne0\)，则
\[
\frac{f(2x-a)}{x-a}
\]
在 \(x\to a^+\) 时分母趋于 \(0^+\)，分子趋于非零常数，极限不可能为有限值。题设说极限存在，故
\[
f(a)=0.
\]
又 \(f'(x)>0\)，所以 \(f\) 在 \([a,b]\) 上严格递增，从而对任意 \(x\in(a,b)\)，有
\[
f(x)>f(a)=0.
\]

\textbf{(2)} 令
\[
F(x)=x^2,\qquad G(x)=\int_a^x f(t)\,dt.
\]
则 \(F,G\) 在 \([a,b]\) 上连续、在 \((a,b)\) 内可导，且
\[
G'(x)=f(x)>0\quad(a<x<b).
\]
由柯西中值定理，存在 \(\xi\in(a,b)\)，使
\[
\frac{F(b)-F(a)}{G(b)-G(a)}=\frac{F'(\xi)}{G'(\xi)}.
\]
即
\[
\frac{b^2-a^2}{\int_a^b f(x)\,dx}
=\frac{2\xi}{f(\xi)}.
\]

\textbf{(3)} 由 \(f(a)=0\)，对区间 \([a,\xi]\) 应用拉格朗日中值定理，存在
\[
\eta\in(a,\xi)
\]
使
\[
f'(\eta)=\frac{f(\xi)-f(a)}{\xi-a}=\frac{f(\xi)}{\xi-a}.
\]
由第 (2) 问
\[
\frac{b^2-a^2}{\int_a^b f(x)\,dx}=\frac{2\xi}{f(\xi)}
\]
可得
\[
\frac{f(\xi)}{\xi-a}(b^2-a^2)
=\frac{2\xi}{\xi-a}\int_a^b f(x)\,dx.
\]
因此
\[
f'(\eta)(b^2-a^2)=\frac{2\xi}{\xi-a}\int_a^b f(x)\,dx.
\]
且 \(\eta\in(a,\xi)\)，所以 \(\eta\ne\xi\)。
\examnote{“极限存在 + 分母趋零”常用于推出端点函数值为零；含两个增量比的结论优先考虑柯西中值定理。}
\end{solutionblock}

\begin{problemblock}
\textbf{49.} 设 \(f(x),g(x)\) 在 \([a,b]\) 上连续，在 \((a,b)\) 内可导，且
\[
g(a)=g(b)=1,\qquad f'(x)\ne0.
\]
试证存在 \(\xi,\eta\in(a,b)\)，使得
\[
\frac{f'(\xi)}{f'(\eta)}
=e^{\xi-\eta}\bigl[g(\xi)+g'(\xi)\bigr].
\]
\end{problemblock}

\begin{solutionblock}
\analysis{目标式中出现 \(e^\xi[g(\xi)+g'(\xi)]\)，这是 \((e^xg(x))'\)。再用一次 \(e^x\) 与 \(f(x)\) 的柯西中值定理，把端点差相同的两个比值连接起来。}
由于 \(f'(x)\ne0\)，由导函数介值性可知 \(f'(x)\) 在 \((a,b)\) 内同号，因此 \(f\) 严格单调，
\[
f(b)-f(a)\ne0.
\]
令
\[
H(x)=e^xg(x).
\]
由 \(g(a)=g(b)=1\)，得
\[
H(b)-H(a)=e^bg(b)-e^ag(a)=e^b-e^a.
\]
对 \(H(x)\) 与 \(f(x)\) 在 \([a,b]\) 上应用柯西中值定理，存在 \(\xi\in(a,b)\)，使
\[
\frac{H(b)-H(a)}{f(b)-f(a)}
=\frac{H'(\xi)}{f'(\xi)}
=\frac{e^\xi[g(\xi)+g'(\xi)]}{f'(\xi)}.
\]
对 \(e^x\) 与 \(f(x)\) 在 \([a,b]\) 上应用柯西中值定理，存在 \(\eta\in(a,b)\)，使
\[
\frac{e^b-e^a}{f(b)-f(a)}
=\frac{e^\eta}{f'(\eta)}.
\]
由于 \(H(b)-H(a)=e^b-e^a\)，两式左端相等，所以
\[
\frac{e^\xi[g(\xi)+g'(\xi)]}{f'(\xi)}
=\frac{e^\eta}{f'(\eta)}.
\]
整理得
\[
\frac{f'(\xi)}{f'(\eta)}
=e^{\xi-\eta}[g(\xi)+g'(\xi)].
\]
\examnote{含 \(g+g'\) 且旁边有指数 \(e^x\) 时，第一反应应是 \((e^xg)'\)。}
\end{solutionblock}

\begin{problemblock}
\textbf{50.} 设函数 \(f(x)\) 在闭区间 \([0,1]\) 上连续，在开区间 \((0,1)\) 内可导，且
\[
f(0)=0,\qquad f(1)=\frac13.
\]
证明：存在
\[
\xi\in\left(0,\frac12\right),\qquad
\eta\in\left(\frac12,1\right),
\]
使得
\[
f'(\xi)+f'(\eta)=\xi^2+\eta^2.
\]
\end{problemblock}

\begin{solutionblock}
\analysis{把右边的 \(\xi^2+\eta^2\) 移到左边，构造 \(F(x)=f(x)-x^3/3\)，使目标变为 \(F'(\xi)+F'(\eta)=0\)。}
令
\[
F(x)=f(x)-\frac{x^3}{3}.
\]
则
\[
F(0)=0,\qquad F(1)=f(1)-\frac13=0.
\]
在区间 \([0,1/2]\) 上应用拉格朗日中值定理，存在
\[
\xi\in\left(0,\frac12\right)
\]
使
\[
F'(\xi)=\frac{F(1/2)-F(0)}{1/2-0}=2F\left(\frac12\right).
\]
在区间 \([1/2,1]\) 上应用拉格朗日中值定理，存在
\[
\eta\in\left(\frac12,1\right)
\]
使
\[
F'(\eta)=\frac{F(1)-F(1/2)}{1-1/2}=-2F\left(\frac12\right).
\]
两式相加得
\[
F'(\xi)+F'(\eta)=0.
\]
而
\[
F'(x)=f'(x)-x^2,
\]
故
\[
f'(\xi)-\xi^2+f'(\eta)-\eta^2=0.
\]
即
\[
f'(\xi)+f'(\eta)=\xi^2+\eta^2.
\]
\examnote{出现 \(x^2\) 作为导数项时，常逆向构造 \(x^3/3\)。左右两个半区间分别用中值定理，是制造“两点导数和”的常用手法。}
\end{solutionblock}

\begin{problemblock}
\textbf{51.} 设 \(f(x)\) 在 \([0,1]\) 上连续，在 \((0,1)\) 内可导，且
\[
f(0)=f(1).
\]
试证存在 \(\xi,\eta\)，满足
\[
0<\xi<\eta<1,
\]
使
\[
f'(\xi)+f'(\eta)=0.
\]
\end{problemblock}

\begin{solutionblock}
\analysis{把区间从中点拆开，两段分别使用拉格朗日中值定理，两个平均斜率正好互为相反数。}
在 \([0,1/2]\) 上应用拉格朗日中值定理，存在
\[
\xi\in\left(0,\frac12\right)
\]
使
\[
f'(\xi)=\frac{f(1/2)-f(0)}{1/2-0}=2\left[f\left(\frac12\right)-f(0)\right].
\]
在 \([1/2,1]\) 上应用拉格朗日中值定理，存在
\[
\eta\in\left(\frac12,1\right)
\]
使
\[
f'(\eta)=\frac{f(1)-f(1/2)}{1-1/2}
=2\left[f(1)-f\left(\frac12\right)\right].
\]
由 \(f(1)=f(0)\)，得
\[
f'(\eta)=2\left[f(0)-f\left(\frac12\right)\right]
=-2\left[f\left(\frac12\right)-f(0)\right].
\]
因此
\[
f'(\xi)+f'(\eta)=0.
\]
又
\[
0<\xi<\frac12<\eta<1,
\]
所以 \(0<\xi<\eta<1\)。
\examnote{两点结论常通过“分段中值定理”获得；中点不是唯一选择，但中点让系数最整齐。}
\end{solutionblock}

\begin{problemblock}
\textbf{52.} 设 \(f(x)\) 在 \([0,1]\) 上连续，在 \((0,1)\) 内可导，且
\[
f(0)=0,\qquad f(1)=0.
\]
若 \(f(x)\) 在 \([0,1]\) 上的最大值为 \(M>0\)，证明存在两个不同的点 \(x_1,x_2\in(0,1)\)，使得
\[
\frac{1}{f'(x_1)}-\frac{1}{f'(x_2)}=\frac{n}{M},
\]
其中 \(n\) 是大于 \(1\) 的整数。
\end{problemblock}

\begin{solutionblock}
\analysis{最大值点左右两侧分别存在正斜率和负斜率。再利用导数的介值性，选出一正一负且大小恰为 \(2M/n\) 的导数值。}
设 \(c\in(0,1)\) 使
\[
f(c)=M.
\]
因为端点值为 \(0\)，且最大值 \(M>0\)，最大值点必在开区间内。由费马定理，
\[
f'(c)=0.
\]
对 \([0,c]\) 应用拉格朗日中值定理，存在 \(u\in(0,c)\)，使
\[
f'(u)=\frac{f(c)-f(0)}{c-0}=\frac{M}{c}\ge M.
\]
因为 \(n>1\) 是整数，所以 \(n\ge2\)，从而
\[
0<\frac{2M}{n}\le M\le f'(u).
\]
导函数具有介值性，且 \(f'(c)=0\)，故在 \((u,c)\) 内存在 \(x_1\)，使
\[
f'(x_1)=\frac{2M}{n}.
\]

同理，对 \([c,1]\) 应用拉格朗日中值定理，存在 \(v\in(c,1)\)，使
\[
f'(v)=\frac{f(1)-f(c)}{1-c}=-\frac{M}{1-c}\le -M.
\]
由于
\[
-M\le-\frac{2M}{n}<0,
\]
并且 \(f'(c)=0,\ f'(v)\le -M\)，由导数介值性，存在 \(x_2\in(c,v)\)，使
\[
f'(x_2)=-\frac{2M}{n}.
\]
于是 \(x_1\ne x_2\)，且
\[
\frac{1}{f'(x_1)}-\frac{1}{f'(x_2)}
=\frac{1}{2M/n}-\frac{1}{-2M/n}
=\frac{n}{2M}+\frac{n}{2M}
=\frac{n}{M}.
\]
\examnote{本题不能只用一次中值定理，还要用导数介值性。导数即使不连续，也一定有介值性，这是考研证明题的高频工具。}
\end{solutionblock}

\begin{problemblock}
\textbf{53.} 设 \(f(x)\) 在 \([0,1]\) 上二阶可导，
\[
f(0)=f(1)=0,\qquad \max_{0\le x\le1}f(x)=2.
\]
试证存在 \(\xi\in(0,1)\)，使得
\[
f''(\xi)\le -16.
\]
\end{problemblock}

\begin{solutionblock}
\analysis{取最大值点 \(c\)，由 \(f(c)=2\) 和 \(f'(c)=0\)，向距离较近的端点作二阶泰勒展开。距离不超过 \(1/2\)，于是二阶导数至少要小到 \(-16\)。}
设 \(c\in(0,1)\) 使
\[
f(c)=2.
\]
由于端点 \(f(0)=f(1)=0\)，最大值点 \(c\) 在开区间内，故
\[
f'(c)=0.
\]
若 \(c\le 1/2\)，对 \(f(x)\) 在点 \(c\) 到点 \(0\) 之间使用泰勒公式的拉格朗日余项，存在
\[
\xi\in(0,c)
\]
使
\[
f(0)=f(c)+f'(c)(0-c)+\frac12 f''(\xi)(0-c)^2.
\]
代入 \(f(0)=0,\ f(c)=2,\ f'(c)=0\)，得
\[
0=2+\frac12 f''(\xi)c^2,
\]
所以
\[
f''(\xi)=-\frac{4}{c^2}.
\]
又 \(c\le1/2\)，故
\[
f''(\xi)=-\frac{4}{c^2}\le -16.
\]

若 \(c>1/2\)，对 \(f(x)\) 在点 \(c\) 到点 \(1\) 之间使用泰勒公式，存在
\[
\xi\in(c,1)
\]
使
\[
f(1)=f(c)+f'(c)(1-c)+\frac12 f''(\xi)(1-c)^2.
\]
同理
\[
0=2+\frac12 f''(\xi)(1-c)^2,
\]
即
\[
f''(\xi)=-\frac{4}{(1-c)^2}.
\]
因为 \(c>1/2\)，所以 \(1-c<1/2\)，从而
\[
f''(\xi)<-16.
\]
综上，必存在 \(\xi\in(0,1)\)，使
\[
f''(\xi)\le -16.
\]
\examnote{最大值为 2 而端点为 0，函数必须“弯下来”。常数 16 来自最近端点距离不超过 \(1/2\)：\(4/(1/2)^2=16\)。}
\end{solutionblock}

\begin{problemblock}
\textbf{54.} 设 \(f(x)\) 在 \([0,2]\) 上二阶可导，且
\[
|f(x)|\le1,\qquad |f''(x)|\le1.
\]
证明：
\[
|f'(x)|\le2\qquad(0\le x\le2).
\]
\end{problemblock}

\begin{solutionblock}
\analysis{反证。若某点导数大于 2，则由于二阶导数绝对值不超过 1，整个区间上的导数都保持较大的同向趋势，积分后会迫使函数值振幅超过 2，与 \(|f|\le1\) 矛盾。}
任取 \(x_0\in[0,2]\)，记
\[
p=f'(x_0).
\]
先证 \(p\le2\)。若 \(p>2\)，由 \(|f''(x)|\le1\)，对任意 \(t\in[0,2]\)，有
\[
f'(t)\ge p-|t-x_0|.
\]
于是
\[
f(2)-f(0)=\int_0^2 f'(t)\,dt
\ge \int_0^2 \bigl(p-|t-x_0|\bigr)\,dt.
\]
计算右端：
\[
\int_0^2 |t-x_0|\,dt
=\frac{x_0^2}{2}+\frac{(2-x_0)^2}{2}\le2.
\]
故
\[
f(2)-f(0)\ge 2p-2>2.
\]
但由 \(|f(x)|\le1\)，有
\[
f(2)-f(0)\le |f(2)|+|f(0)|\le2,
\]
矛盾。因此
\[
f'(x_0)\le2.
\]

再证 \(p\ge-2\)。若 \(p<-2\)，同理对任意 \(t\in[0,2]\)，
\[
f'(t)\le p+|t-x_0|.
\]
于是
\[
f(2)-f(0)=\int_0^2 f'(t)\,dt
\le \int_0^2 \bigl(p+|t-x_0|\bigr)\,dt
\le 2p+2<-2.
\]
但
\[
f(2)-f(0)\ge -|f(2)|-|f(0)|\ge-2,
\]
矛盾。因此
\[
f'(x_0)\ge-2.
\]
由于 \(x_0\in[0,2]\) 任意，故
\[
|f'(x)|\le2\qquad(0\le x\le2).
\]
\examnote{这类题不要只用单点泰勒估计；把导数在全区间积分，能同时利用 \(|f|\le1\) 与 \(|f''|\le1\)。}
\end{solutionblock}

\section{第二章完成说明}
第二章第 \(1\)--\(54\) 题已按“题目解析、完整推导、答案、考研提示”的格式整理完毕。下一步从第三章一元函数积分学继续补充。
